\scanpage{79}
On est obligé de renforcer la propriété universelle
de $\otimes$ et $\Gamma^p$, pour avoir les propriétés d'assoc. etc
qu'on veut.

Soit $X\times Y\longrightarrow\deleted{\unsure{P}}\deleted{\unsure{P}}T$ un morphisme de bidegré $(1,1)$,
alors pour tt $Z$ on a
\[
X\times Y\times Z\longrightarrow T\times Z
\]
d'où pour tt $M$
\[
P_{1,n}(T,Z;M)\longrightarrow P_{1,1,n}(X,Y,Z;M)
\]
On dit que $T$ est un produit \uline{tensoriel strict} de
$X$ et $Y$ si l'hom précédent est un iso,
\deleted{et \unsure{alors} $X\otimes Y=T$}
pour tt $n$ et tout $M$. Faisant $n=\unsure{1|0}$, \inserted{$Z=0$}, on trouve
\[
P_{1,\deleted{\unsure{1}}}(T,M)\simeq P_{1,1}(X,Y;M)
\]
donc $T$ est un produit tensoriel. C'est
en exigeant qu'on a des $\otimes$ \uline{stricts} \deleted{pour}\inserted{\unsure{de deux}} et
trois facteurs (\deleted{\unsure{qu'il}} il faut que la définition précédente
\inserted{au cas d'un produit \unsure{tensoriel de $n$ facteurs}}) que l'on
trouve \unsure{automatiqt} une structure
ass. comm. compatible avec
\[
\deleted{\operatorname{Hom}_{P_1}} P_1(X_1\otimes\cdots\otimes X_r;M)\simeq P_{1,\ldots,1}(X_1,\ldots,X_r;M)
\]
On sait qu'un objet unité (sous réserve d'une
donnée d'associativité) si elle existe, est unique.
\margin{?} Il faut examiner pourquoi ce devrait être $\Gamma^0(0)$.

Revenons aux $\Gamma^n$, soit $X=X_1\times\cdots\times X_r$,
tel que $\Gamma^n(X)$ existe. Considérons les
sous foncteurs $P_{\nu=(\nu_1,\ldots,\nu_r)}(X_1,\ldots,X_r;M)$
du foncteur \inserted{$P_1\longrightarrow(\mathrm{Ab})$} (cov.) représenté par $\Gamma^n(X)$,
supposons qu'ils soient aussi \unsure{corepr.|repr.} (automatique
si $P_1$ est karoubienne) soit $\Gamma^{\nu_1,\ldots,\nu_r}(X_1,\ldots,X_r)$
$=\Gamma^\nu(X_1,\ldots,X_r)$ l'\deleted{\unsure{foncteur}}\inserted{objet} qui le représente. $\deleted{\unsure{X_1}}$
On va définir \unsure{leur notion|leurs niveaux|sans réserve|sous réserve} d'existence des $\Gamma^{\nu_i}X_i$
\deleted{produit}\inserted{et des} tensoriel, \deleted{\unsure{en}} ou dans $\check{P}_1=\widehat{P_1^\circ}^{\,\circ}$ si.
% archive sheet number: 78
